% Examenes_de_PF_con_Haskell_Vol11.tex
% Exámenes de programación funcional con Haskell (Curso 2019-20)
% José A. Alonso Jiménez <jalonso@us.es>
% Sevilla, 31 de octubre de 2019
% =============================================================================

\documentclass[a4paper,12pt,twoside]{book}

%%%%%%%%%%%%%%%%%%%%%%%%%%%%%%%%%%%%%%%%%%%%%%%%%%%%%%%%%%%%%%%%%%%%%%%%%%%%%%
%% § Paquetes adicionales                                                   %%
%%%%%%%%%%%%%%%%%%%%%%%%%%%%%%%%%%%%%%%%%%%%%%%%%%%%%%%%%%%%%%%%%%%%%%%%%%%%%%

% Configuración para XeLaTeX
\usepackage{fontspec}
\usepackage{xltxtra}
\defaultfontfeatures{Ligatures=TeX,Numbers=OldStyle}
\setromanfont{DejaVu Sans}
% \setsansfont{Arial}
\setmonofont{DejaVu Sans Mono}[Scale={0.90}]

% Notas: La lista de fuentes disponibles se obtiene con fc-list

% \usepackage{ucs}
% \usepackage[utf8]{inputenc}        % Acentos de UTF8
\usepackage[spanish]{babel}        % Castellanización.
% \usepackage[T1]{fontenc}           % Codificación T1 con European Computer
%                                    % Modern.  
% \usepackage{graphicx}
\usepackage{fancyvrb}              % Verbatim extendido
% \usepackage{mathpazo}              % Fuentes semejante a palatino
% \usepackage[scaled=.90]{helvet}
% \usepackage{cmtt}
% \renewcommand{\ttdefault}{cmtt}
\usepackage{a4wide}
\usepackage{minted}

\linespread{1.05}                  % Distancia entre líneas
\setlength{\parindent}{2em}        % Indentación de comienzo de párrafo
% \deactivatetilden                  % Elima uso de ~ para la eñe
\raggedbottom                      % No ajusta los espacios verticales.

\usepackage[%
  colorlinks=true,
  urlcolor=blue,
  % pdftex,
  pdfauthor={José A. Alonso <jalonso@us.es>},%
  pdftitle={Examenes de Informatica de 1 de Matematicas},%
  pdfstartview=FitH,%
  bookmarks=false]{hyperref}      

\setcounter{tocdepth}{1}
\setcounter{secnumdepth}{4}

\usepackage{tocstyle}
\usetocstyle{KOMAlike}

\newenvironment{enumerate*}%
  {\vspace*{-0mm}
   \begin{enumerate}%
    \setlength{\itemsep}{0pt}%
    \setlength{\parskip}{0pt}}%
  {\vspace*{-0mm}
   \end{enumerate}}

\newenvironment{itemize*}%
  {\vspace*{-0mm}
   \begin{itemize}%
    \setlength{\itemsep}{0pt}%
    \setlength{\parskip}{0pt}}%
  {\vspace*{-0mm}
   \end{itemize}}

%%%%%%%%%%%%%%%%%%%%%%%%%%%%%%%%%%%%%%%%%%%%%%%%%%%%%%%%%%%%%%%%%%%%%%%%%%%%%%
%% § Cabeceras                                                              %%
%%%%%%%%%%%%%%%%%%%%%%%%%%%%%%%%%%%%%%%%%%%%%%%%%%%%%%%%%%%%%%%%%%%%%%%%%%%%%%

\usepackage{fancyhdr}

\addtolength{\headheight}{\baselineskip}

\pagestyle{fancy}

\cfoot{}

\fancyhead{}
\fancyhead[RE]{\mdseries\sffamily \nouppercase{\leftmark}}
\fancyhead[LO]{\mdseries\sffamily \nouppercase{\rightmark}}
\fancyhead[LE,RO]{\mdseries\sffamily \thepage}

%%%%%%%%%%%%%%%%%%%%%%%%%%%%%%%%%%%%%%%%%%%%%%%%%%%%%%%%%%%%%%%%%%%%%%%%%%%%%%
%% § Definiciones                                                           %%
%%%%%%%%%%%%%%%%%%%%%%%%%%%%%%%%%%%%%%%%%%%%%%%%%%%%%%%%%%%%%%%%%%%%%%%%%%%%%%

\input definiciones

%%%%%%%%%%%%%%%%%%%%%%%%%%%%%%%%%%%%%%%%%%%%%%%%%%%%%%%%%%%%%%%%%%%%%%%%%%%%%%
%% § Título                                                                 %%
%%%%%%%%%%%%%%%%%%%%%%%%%%%%%%%%%%%%%%%%%%%%%%%%%%%%%%%%%%%%%%%%%%%%%%%%%%%%%%

\title{
  {\LARGE Exámenes de \\ ``Programaci\'on funcional con Haskell''} \\ 
  {\large Vol. 11 (Curso 2019--20)}} 
\author{
  \href{http://www.cs.us.es/~jalonso}{José A. Alonso Jiménez}}
\date{\vfill \hrule \vspace*{2mm}
  \begin{tabular}{l}
      \href{http://www.cs.us.es/glc}
           {Grupo de Lógica Computacional} \\
      \href{http://www.cs.us.es}
           {Dpto. de Ciencias de la Computación e Inteligencia Artificial} \\
      \href{http://www.us.es}
           {Universidad de Sevilla}  \\
      Sevilla, 31 de octubre de 2019
  \end{tabular}\hfill\mbox{}}

%%%%%%%%%%%%%%%%%%%%%%%%%%%%%%%%%%%%%%%%%%%%%%%%%%%%%%%%%%%%%%%%%%%%%%%%%%%%%%
%% § Documento                                                              %%
%%%%%%%%%%%%%%%%%%%%%%%%%%%%%%%%%%%%%%%%%%%%%%%%%%%%%%%%%%%%%%%%%%%%%%%%%%%%%%

% \includeonly{}

% \includexmp{licencia}

\begin{document}

\maketitle
\newpage

\noindent
Esta obra está bajo una licencia Reconocimiento--NoComercial--CompartirIgual
2.5 Spain de Creative Commons. 

\vspace*{4ex}

\begin{center}
\fbox{
\begin{minipage}{\textwidth}
\vspace*{1ex}
\noindent                     
{\bf Se permite:}
\begin{itemize}                         
\item copiar, distribuir y comunicar públicamente la obra
\item hacer obras derivadas
\end{itemize}
         
\noindent       
{\bf Bajo las condiciones siguientes:}

\begin{minipage}[t]{20mm}
  \vspace*{0mm}\hspace*{1em}
  \includegraphics[scale=0.35]{Licencia/deed.png}
\end{minipage}%
\begin{minipage}[t]{14cm}
  \vspace*{0mm}
  {\bf Reconocimiento}. Debe reconocer los créditos de la  obra de la manera
  especificada por el autor.  
\end{minipage}
\\[1ex]

\begin{minipage}[t]{20mm}
  \vspace*{0mm}\hspace*{1em}
  \includegraphics[scale=0.35]{Licencia/deed-eu.png}
\end{minipage}%
\begin{minipage}[t]{14cm}
  \vspace*{2mm}
  {\bf No comercial}. No puede utilizar esta obra para fines comerciales.
\end{minipage}
\\[1ex]

\begin{minipage}[t]{20mm}
  \vspace*{0mm}\hspace*{1em}
  \includegraphics[scale=0.35]{Licencia/deed_002.png}
\end{minipage}%
\begin{minipage}[t]{14cm}
  \vspace*{-2mm}
  {\bf Compartir bajo la misma licencia}. Si altera o transforma esta obra, o
  genera una obra derivada, sólo puede distribuir la obra generada bajo una
  licencia idéntica a ésta. 
\end{minipage}

\begin{itemize}    
\item Al reutilizar o distribuir la obra, tiene que dejar bien claro los
  términos de la licencia de esta obra. 
\item Alguna de estas condiciones puede no aplicarse si se obtiene el permiso
 del titular de los derechos de autor.
\end{itemize}
\vspace*{1ex}
\end{minipage}}
\end{center}

\vspace*{4ex}

Esto es un resumen del texto legal (la licencia completa). Para ver una copia
de esta licencia, visite 
\href{http://creativecommons.org/licenses/by-nc-sa/2.5/es/}
     {\url{http://creativecommons.org/licenses/by-nc-sa/2.5/es/}}
o envie una carta a Creative Commons, 559 Nathan Abbott Way, Stanford,
California 94305, USA.

%%% Local Variables: 
%%% mode: latex
%%% TeX-master: t
%%% End: 

\newpage

\tableofcontents
\clearpage

\renewcommand{\chaptername}{}

\chapter*{Introducción}
\addcontentsline{toc}{chapter}{Introducción}

Este libro es una recopilación de las soluciones de
ejercicios de los exámenes de programación funcional con Haskell de la
\href{http://www.cs.us.es/~jalonso/cursos/i1m-19}
     {asignatura de Informática (curso 2019--20)}
del
\href{http://www.matematicas.us.es/estudios/grado-en-matematicas}
     {Grado en Matemática} 
de la 
\href{http://www.us.es/}
     {Universidad de Sevilla}.

Los exámenes se realizaron en el aula de informática y su duración
fue de 2 horas. La materia de cada examen es la impartida desde el
comienzo del curso (generalmente, el 1 de octubre) hasta la fecha
del examen. Dicha materia se encuentra en los libros de temas y
ejercicios del curso:
\begin{itemize}
\item
  \href{https://www.cs.us.es/~jalonso/cursos/i1m-19/temas/2019-20-IM-temas-PF.pdf}
  {Temas de programación funcional (curso 2019--20)}\
  \footnote{\url{https://www.cs.us.es/~jalonso/cursos/i1m-19/temas/2019-20-IM-temas-PF.pdf}} 
\item
  \href{https://www.cs.us.es/~jalonso/cursos/i1m-19/ejercicios/ejercicios-I1M-2019.pdf}
  {Ejercicios de ``Informática de 1º de Matemáticas'' (2019--20)}\
  \footnote{\url{https://www.cs.us.es/~jalonso/cursos/i1m-19/ejercicios/ejercicios-I1M-2019.pdf}}
\item
  \href{http://www.cs.us.es/~jalonso/publicaciones/Piensa_en_Haskell.pdf}
  {Piensa en Haskell (Ejercicios de programación funcional con Haskell)}\
  \footnote{\url{http://www.cs.us.es/~jalonso/publicaciones/Piensa_en_Haskell.pdf}}
\end{itemize}

El libro consta de 5 capítulos correspondientes a 5 grupos de la
asignatura. En cada capítulo hay una sección por cada uno de los
exámenes del grupo. Los ejercicios de cada examen han sido propuestos
por los profesores de su grupo (cuyos nombres aparecen en el título del
capítulo). Sin embargo, los he modificado para unificar el estilo de su
presentación.

Finalmente, el libro contiene dos apéndices. Uno con el método de Polya
de resolución de problemas (sobre el que se hace énfasis durante todo el
curso) y el otro con un resumen de las funciones de Haskell de uso más
frecuente.

Los códigos del libro están disponibles en
\href{https://github.com/jaalonso/Examenes_de_PF_con_Haskell_Vol11}
     {GitHub}
     \footnote{{\url{https://github.com/jaalonso/Examenes_de_PF_con_Haskell_Vol11}}}

Este libro es el 11º volumen de la serie de recopilaciones de
exámenes de programación funcional con Haskell. Los volúmenes anteriores son
\begin{itemize}
\item
  \href{https://github.com/jaalonso/Examenes_de_PF_con_Haskell_Vol1}
  {Exámenes de ``Programaci\'on funcional con Haskell''.
    Vol. 1 (Curso 2009--10)}\
    \footnote{\url{{https://github.com/jaalonso/Examenes_de_PF_con_Haskell_Vol1}}}

\item
  \href{https://github.com/jaalonso/Examenes_de_PF_con_Haskell_Vol2}
  {Exámenes de ``Programaci\'on funcional con Haskell''.
    Vol. 2 (Curso 2010--11)}\
    \footnote{\url{{https://github.com/jaalonso/Examenes_de_PF_con_Haskell_Vol2}}}

\item
  \href{https://github.com/jaalonso/Examenes_de_PF_con_Haskell_Vol3}
  {Exámenes de ``Programaci\'on funcional con Haskell''.
    Vol. 3 (Curso 2011--12)}\
    \footnote{\url{{https://github.com/jaalonso/Examenes_de_PF_con_Haskell_Vol3}}}

\item
  \href{https://github.com/jaalonso/Examenes_de_PF_con_Haskell_Vol4}
  {Exámenes de ``Programaci\'on funcional con Haskell''.
    Vol. 4 (Curso 2012--13)}\
    \footnote{\url{{https://github.com/jaalonso/Examenes_de_PF_con_Haskell_Vol4}}}

\item
  \href{https://github.com/jaalonso/Examenes_de_PF_con_Haskell_Vol5}
  {Exámenes de ``Programaci\'on funcional con Haskell''.
    Vol. 5 (Curso 2013--14)}\
    \footnote{\url{{https://github.com/jaalonso/Examenes_de_PF_con_Haskell_Vol5}}}

\item
  \href{https://github.com/jaalonso/Examenes_de_PF_con_Haskell_Vol6}
  {Exámenes de ``Programaci\'on funcional con Haskell''.
    Vol. 6 (Curso 2014--15)}\
    \footnote{\url{{https://github.com/jaalonso/Examenes_de_PF_con_Haskell_Vol6}}}

\item
  \href{https://github.com/jaalonso/Examenes_de_PF_con_Haskell_Vol7}
  {Exámenes de ``Programaci\'on funcional con Haskell''.
    Vol. 6 (Curso 2015--16)}\
    \footnote{\url{{https://github.com/jaalonso/Examenes_de_PF_con_Haskell_Vol7}}}

\item
  \href{https://github.com/jaalonso/Examenes_de_PF_con_Haskell_Vol8}
  {Exámenes de ``Programaci\'on funcional con Haskell''.
    Vol. 8 (Curso 2016--17)}\
    \footnote{\url{{https://github.com/jaalonso/Examenes_de_PF_con_Haskell_Vol8}}}

\item
  \href{https://github.com/jaalonso/Examenes_de_PF_con_Haskell_Vol9}
  {Exámenes de ``Programaci\'on funcional con Haskell''.
    Vol. 9 (Curso 2017--18)}\
    \footnote{\url{{https://github.com/jaalonso/Examenes_de_PF_con_Haskell_Vol9}}}

\item
  \href{https://github.com/jaalonso/Examenes_de_PF_con_Haskell_Vol10}
  {Exámenes de ``Programaci\'on funcional con Haskell''.
    Vol. 10 (Curso 2018--19)}\
    \footnote{\url{{https://github.com/jaalonso/Examenes_de_PF_con_Haskell_Vol10}}}

\end{itemize}
     
\begin{flushright}
  José A. Alonso \\
  Sevilla, 31 de octubre de 2019
\end{flushright}

% \chapter{Exámenes del grupo 1}
% \chapterauthor{Antonia M. Chávez}
% \section{Examen 1 (11 de noviembre de 2019)}
% \examen{../Examenes/Grupo_1/examen_1_11_nov.hs}
% \section{Examen 2 (16 de diciembre de 2019)}
% \examen{../Examenes/Grupo_1/examen_12_16_dic.hs}
% \section{Examen 3 (23 de enero de 2020)}
% \examen{../Examenes/Grupo_1/examen_3_23_ene.hs}
% \section{Examen 4 (16 de marzo de 2020}
% \examen{../Examenes/Grupo_1/examen_4_16_mar.hs}
% \section{Examen 5 (11 de mayo de 2020)}
% \examen{../Examenes/Grupo_1/examen_5_11_may.hs}
% \section{Examen 6 (15 de junio de 2020)}
% \examen{../Examenes/Grupo_1/examen_6_15_jun.hs}
% \section{Examen 7 (23 de junio de 2020)}
% \examen{../Examenes/Grupo_1/examen_7_23_jun.hs}
% \section{Examen 8 (02 de septiembre de 2020)}
% \examen{../Examenes/Grupo_1/examen_8_02_sep.hs}
        
% \chapter{Exámenes del grupo 2}
% \chapterauthor{María J. Hidalgo}
% \section{Examen 1 (06 de noviembre de 2019)}
% \examen{../Examenes/Grupo_2/examen_1_06_nov.hs}
% \section{Examen 2 (27 de noviembre de 2019)}
% \examen{../Examenes/Grupo_2/examen_2_27_nov.hs}
% \section{Examen 3 (23 de enero de 2020)}
% \examen{../Examenes/Grupo_2/examen_3_23_ene.hs}
% \section{Examen 4 (13 de marzo de 2020)}
% \examen{../Examenes/Grupo_2/examen_4_13_mar.hs}
% \section{Examen 5 (24 de abril de 2020)}
% \examen{../Examenes/Grupo_2/examen_5_24_abr.hs}
% \section{Examen 6 (15 de junio de 2020)}
% \examen{../Examenes/Grupo_2/examen_6_15_jun.hs}
% \section{Examen 7 (23 de junio de 2020)}
% \examen{../Examenes/Grupo_2/examen_7_23_jun.hs}
% \section{Examen 8 (02 de septiembre de 2020)}
% \examen{../Examenes/Grupo_2/examen_8_02_sep.hs}
    
\chapter{Exámenes del grupo 3}
\chapterauthor{Francisco J. Martín}
\section{Examen 1 (28 de octubre de 2019}
\examen{../Examenes/Grupo_3/examen_1_28_oct.hs}
% \section{Examen 2 (02 de diciembre de 2019}
% \examen{../Examenes/Grupo_3/examen_2_02_dic.hs}
% \section{Examen 3 (23 de enero de 2020)}
% \examen{../Examenes/Grupo_3/examen_3_23_ene.hs}
% \section{Examen 4 (16 de marzo de 2020}
% \examen{../Examenes/Grupo_3/examen_4_16_mar.hs}
% \section{Examen 5 (04 de mayo de 2020}
% \examen{../Examenes/Grupo_3/examen_5_04_may.hs}
% \section{Examen 6 (15 de junio de 2020)}
% \examen{../Examenes/Grupo_3/examen_6_15_jun.hs}
% \section{Examen 7 (23 de junio de 2020)}
% \examen{../Examenes/Grupo_3/examen_7_23_jun.hs}
% \section{Examen 8 (02 de septiembre de 2020)}
% \examen{../Examenes/Grupo_3/examen_8_02_sep.hs}

\chapter{Exámenes del grupo 4}
\chapterauthor{José A. Alonso}
\section{Examen 1 (30 de octubre de 2019)}
\examen{../Examenes/Grupo_4/examen_1_30_oct.hs}
% \section{Examen 2 (18 de diciembre de 2019)}
% \examen{../Examenes/Grupo_4/examen_2_18_dic.hs}
% \section{Examen 3 (23 de enero de 2020)}
% \examen{../Examenes/Grupo_4/examen_3_23_ene.hs}
% \section{Examen 4 (18 de marzo de 2020)}
% \examen{../Examenes/Grupo_4/examen_4_18_mar.hs}
% \section{Examen 5 (06 de mayo de 2020)}
% \examen{../Examenes/Grupo_4/examen_5_06_may.hs}
% \section{Examen 6 (15 de junio de 2020)}
% \examen{../Examenes/Grupo_4/examen_6_15_jun.hs}
% \section{Examen 7 (23 de junio de 2020)}
% \examen{../Examenes/Grupo_4/examen_7_23_jun.hs}
% \section{Examen 8 (02 de septiembre de 2020)}
% \examen{../Examenes/Grupo_4/examen_8_02_sep.hs}

% \chapter{Exámenes del grupo 5}
% \chapterauthor{Andrés Cordón y Miguel A. Martínez}
% \section{Examen 1 (06 de noviembre de 2019)}
% \examen{../Examenes/Grupo_5/examen_1_06_nov.hs}
% \section{Examen 2 (11 de diciembre de 2019)}
% \examen{../Examenes/Grupo_5/examen_2_11_dic.hs}
% \section{Examen 3 (23 de enero de 2020)}
% \examen{../Examenes/Grupo_5/examen_3_23_ene.hs}
% \section{Examen 4 (11 de marzo de 2020)}
% \examen{../Examenes/Grupo_5/examen_4_11_mar.hs}
% \section{Examen 5 (06 de mayo de 2020)}
% \examen{../Examenes/Grupo_5/examen_5_06_may.hs}
% \section{Examen 6 (15 de junio de 2020)}
% \examen{../Examenes/Grupo_5/examen_6_15_jun.hs}
% \section{Examen 7 (23 de junio de 2020)}
% \examen{../Examenes/Grupo_5/examen_7_23_jun.hs}
% \section{Examen 8 (02 de septiembre de 2020)}
% \examen{../Examenes/Grupo_5/examen_8_02_sep.hs}

\appendix % Apéndices

% A. Resumen de funciones Haskell
% resumen_Haskell.tex
% Resumen de funciones habituales de Haskell.
% José A. Alonso Jiménez <jalonso@us,es>
% Sevilla, 13 de Noviembre de 2007
% ============================================================================

%%%%%%%%%%%%%%%%%%%%%%%%%%%%%%%%%%%%%%%%%%%%%%%%%%%%%%%%%%%%%%%%%%%%%%%%%%%%%%
%% § Paquetes adicionales                                                   %%
%%%%%%%%%%%%%%%%%%%%%%%%%%%%%%%%%%%%%%%%%%%%%%%%%%%%%%%%%%%%%%%%%%%%%%%%%%%%%%

\newcommand{\verba}[1]{%
  \fbox{\textcolor{blue}{\ \texttt{#1}\phantom{I}}}}

%%%%%%%%%%%%%%%%%%%%%%%%%%%%%%%%%%%%%%%%%%%%%%%%%%%%%%%%%%%%%%%%%%%%%%%%%%%%%%
%% § Documento                                                              %%
%%%%%%%%%%%%%%%%%%%%%%%%%%%%%%%%%%%%%%%%%%%%%%%%%%%%%%%%%%%%%%%%%%%%%%%%%%%%%%

\chapter{Resumen de funciones predefinidas de Haskell}
\label{resumen}

\begin{enumerate*}
\item \verba{x + y} es la suma de \verb|x| e \verb|y|.
\item \verba{x - y} es la resta de \verb|x| e \verb|y|.
\item \verba{x / y} es el cociente de \verb|x| entre \verb|y|.
\item \verba{x \^\ y} es \verb|x| elevado a \verb|y|.
\item \verba{x == y} se verifica si \verb|x| es igual a \verb|y|.
\item \verba{x /= y} se verifica si \verb|x| es distinto de \verb|y|.
\item \verba{x <\ y} se verifica si \verb|x| es menor que \verb|y|.
\item \verba{x <= y} se verifica si \verb|x| es menor o igual que \verb|y|.
\item \verba{x >\ y} se verifica si \verb|x| es mayor que \verb|y|.
\item \verba{x >= y} se verifica si \verb|x| es mayor o igual que \verb|y|.
\item \verba{x \&\& y} es la conjunción de \verb|x| e \verb|y|.
\item \verba{x || y} es la disyunción de \verb|x| e \verb|y|.
\item \verba{x:ys} es la lista obtenida añadiendo \verb|x| al principio de
  \verb|ys|. 
\item \verba{xs ++ ys} es la concatenación de \verb|xs| e \verb|ys|.
\item \verba{xs !! n} es el elemento \verb|n|--ésimo de \verb|xs|.
\item \verba{f . g} es la composición de \verb|f| y \verb|g|.
\item \verba{abs x} es el valor absoluto de \verb|x|.
\item \verba{and xs} es la conjunción de la lista de booleanos \verb|xs|.
\item \verba{ceiling x} es el menor entero no menor que \verb|x|.
\item \verba{chr n} es el carácter cuyo código ASCII es \verb|n|.
\item \verba{concat xss} es la concatenación de la lista de listas \verb|xss|.
\item \verba{const x y} es \verb|x|.
\item \verba{curry f} es la versión curryficada de la función \verb|f|.
\item \verba{div x y} es la división entera de \verb|x| entre \verb|y|.
\item \verba{drop n xs} borra los \verb|n| primeros elementos de \verb|xs|.
\item \verba{dropWhile p xs} borra el mayor prefijo de \verb|xs| cuyos
  elementos satisfacen el predicado \verb|p|.
\item \verba{elem x ys} se verifica si \verb|x| pertenece a \verb|ys|.
\item \verba{even x} se verifica si \verb|x| es par.
\item \verba{filter p xs} es la lista de elementos de la lista \verb|xs| que
  verifican el predicado \verb|p|. 
\item \verba{flip f x y} es \verb|f y x|.
\item \verba{floor x} es el mayor entero no mayor que \verb|x|.
\item \verba{foldl f e xs} pliega \verb|xs| de izquierda a derecha
  usando el operador f y el valor inicial \verb|e|.
\item \verba{foldr f e xs} pliega \verb|xs| de derecha a izquierda
  usando el operador f y el valor inicial \verb|e|. 
\item \verba{fromIntegral x} transforma el número entero \verb|x| al tipo
  numérico correspondiente. 
\item \verba{fst p} es el primer elemento del par \verb|p|.
\item \verba{gcd x y} es el máximo común divisor de de \verb|x| e \verb|y|.
\item \verba{head xs} es el primer elemento de la lista \verb|xs|.
\item \verba{init xs} es la lista obtenida eliminando el último elemento de
  \verb|xs|. 
% \item \verba{isSpace x}    se verifica si \verb|x| es un espacio.
% \item \verba{isUpper x}    se verifica si \verb|x| está en mayúscula.
% \item \verba{isLower x}    se verifica si \verb|x| está en minúscula.
% \item \verba{isAlpha x}    se verifica si \verb|x| es un carácter alfabético.
% \item \verba{isDigit x}    se verifica si \verb|x| es un dígito.
% \item \verba{isAlphaNum x} se verifica si \verb|x| es un carácter alfanumérico.
\item \verba{iterate f x} es la lista \verb|[x, f(x), f(f(x)), ...]|.
\item \verba{last xs} es el último elemento de la lista \verb|xs|.
\item \verba{length xs} es el número de elementos de la lista \verb|xs|.
\item \verba{map f xs} es la lista obtenida aplicado \verb|f| a cada elemento
  de \verb|xs|.
\item \verba{max x y} es el máximo de \verb|x| e \verb|y|.
\item \verba{maximum xs} es el máximo elemento de la lista \verb|xs|.
\item \verba{min x y} es el mínimo de \verb|x| e \verb|y|.
\item \verba{minimum xs} es el mínimo elemento de la lista \verb|xs|.
\item \verba{mod x y} es el resto de \verb|x| entre \verb|y|.
\item \verba{not x} es la negación lógica del booleano \verb|x|.
\item \verba{noElem x ys} se verifica si \verb|x| no pertenece a \verb|ys|.
\item \verba{null xs} se verifica si \verb|xs| es la lista vacía.
\item \verba{odd x} se verifica si \verb|x| es impar.
\item \verba{or xs} es la disyunción de la lista de booleanos \verb|xs|.
\item \verba{ord c} es el código ASCII del carácter \verb|c|.
\item \verba{product xs} es el producto de la lista de números \verb|xs|.
\item \verba{read c} es la expresión representada por la cadena \verb|c|.
\item \verba{rem x y} es el resto de \verb|x| entre \verb|y|.
\item \verba{repeat x} es la lista infinita \verb|[x, x, x, ...]|.
\item \verba{replicate n x} es la lista formada por \verb|n| veces el elemento
  \verb|x|. 
\item \verba{reverse xs} es la inversa de la lista \verb|xs|.
\item \verba{round x} es el redondeo de \verb|x| al entero más cercano.
\item \verba{scanr f e xs} es la lista de los resultados de plegar \verb|xs|
  por la derecha con \verb|f| y \verb|e|.
\item \verba{show x} es la representación de \verb|x| como cadena.
\item \verba{signum x} es 1 si \verb|x| es positivo, 0 si \verb|x| es cero y -1
  si \verb|x| es negativo.
\item \verba{snd p} es el segundo elemento del par \verb|p|.
\item \verba{splitAt n xs} es \verb|(take n xs, drop n xs)|.
\item \verba{sqrt x} es la raíz cuadrada de \verb|x|.
\item \verba{sum xs} es la suma de la lista numérica \verb|xs|.
\item \verba{tail xs} es la lista obtenida eliminando el primer elemento de
  \verb|xs|. 
\item \verba{take n xs} es la lista de los \verb|n| primeros elementos de
  \verb|xs|. 
\item \verba{takeWhile p xs} es el mayor prefijo de \verb|xs| cuyos elementos
  satisfacen el predicado \verb|p|. 
\item \verba{uncurry f} es la versión cartesiana de la función \verb|f|.
\item \verba{until p f x} aplica \verb|f| a \verb|x| hasta que se verifique
  \verb|p|. 
\item \verba{zip xs ys} es la lista de pares formado por los correspondientes
  elementos de \verb|xs| e \verb|ys|.
\item \verba{zipWith f xs ys} se obtiene aplicando \verb|f| a los
  correspondientes elementos de \verb|xs| e \verb|ys|. 
\end{enumerate*}

\section{Resumen de funciones sobre TAD en Haskell}

\subsection{Polinomios}

\begin{enumerate*}
\item \verba{polCero} es el polinomio cero.
\item \verba{(esPolCero p)} se verifica si \verb|p| es el polinomio cero.
\item \verba{(consPol n b p)} es el polinomio $bx^n+p$.
\item \verba{(grado p)} es el grado del polinomio \verb|p|.
\item \verba{(coefLider p)} es el coeficiente líder del polinomio \verb|p|.
\item \verba{(restoPol p)} es el resto del polinomio \verb|p|.
\end{enumerate*}

\subsection{Vectores y matrices (\texttt{Data.Array})}

\begin{enumerate*}
\item \verba{(range m n)} es la lista de los índices del \verb|m| al \verb|n|.
\item \verba{(index (m,n) i)} es el ordinal del índice \verb|i| en \verb|(m,n)|.
\item \verba{(inRange (m,n) i)} se verifica si el índice \verb|i| está dentro
  del rango limitado por \verb|m| y \verb|n|.
\item \verba{(rangeSize (m,n))} es el número de elementos en el rango
  limitado por \verb|m| y \verb|n|.
\item \verba{(array (1,n) [(i, f i) | i <- [1..n])} es el vector de dimensión
  \verb|n| cuyo elemento \verb|i|--ésimo es \verb|f i|. 
\item \verba{(array ((1,1),(m,n)) [((i,j), f i j) | i <- [1..m], j <- [1..n]])} 
  es la matriz de dimensión \verb|m.n| cuyo elemento \verb|(i,j)|--ésimo es
  \verb|f i j|.  
\item \verba{(array (m,n) ivs)} es la tabla de índices en el rango limitado
  por \verb|m| y \verb|n| definida por la lista de asociación \verb|ivs|
  (cuyos elementos son pares de la forma (índice, valor)).
\item \verba{(t ! i)} es el valor del índice \verb|i| en la tabla \verb|t|.
\item \verba{(bounds t)} es el rango de la tabla \verb|t|.
\item \verba{(indices t)} es la lista de los índices de la tabla \verb|t|.
\item \verba{(elems t)} es la lista de los elementos de la tabla \verb|t|.
\item \verba{(assocs t)} es la lista de asociaciones de la tabla \verb|t|.
\item \verba{(t // ivs)} es la tabla \verb|t| asignándole a los índices de la
  lista de asociación \verb|ivs| sus correspondientes valores.
\item \verba{(listArray (m,n) vs)} es la tabla cuyo rango es \verb|(m,n)| y
  cuya lista de valores es \verb|vs|.
\item \verba{(accumArray f v (m,n) ivs)} es la tabla de rango \verb|(m,n)|
  tal que el valor del índice \verb|i| se obtiene acumulando la aplicación de
  la función \verb|f| al valor inicial \verb|v| y a los valores de la lista
  de asociación \verb|ivs| cuyo índice es \verb|i|.
\end{enumerate*}

\subsection{Tablas}
\begin{enumerate*}
\item \verba{(tabla ivs)} es la tabla correspondiente a la lista de
  asociación \verb|ivs| (que es una lista de pares formados por los
  índices y los valores).
\item \verba{(valor t i)} es el valor del índice \verb|i| en la tabla
  \verb|t|.
\item \verba{(modifica (i,v) t)} es la tabla obtenida modificando en la
  tabla \verb|t| el valor de \verb|i| por \verb|v|.
\end{enumerate*}

\subsection{Grafos}
\begin{enumerate*}
  \item \verba{(creaGrafo d cs as)} es un grafo (dirigido o no, según el valor
    de o), con el par de cotas cs y listas de aristas as (cada arista es un
    trío formado por los dos vértices y su peso).
  \item \verba{(dirigido g)} se verifica si \verb|g| es dirigido.
  \item \verba{(nodos g)} es la lista de todos los nodos del grafo \verb|g|.
  \item \verba{(aristas g)} es la lista de las aristas del grafo \verb|g|.
  \item \verba{(adyacentes g v)} es la lista de los vértices adyacentes al nodo
    \verb|v| en el grafo \verb|g|.
  \item \verba{(aristaEn g a)} se verifica si \verb|a| es una arista del grafo
    \verb|g|.
  \item \verba{(peso v1 v2 g)} es el peso de la arista que une los vértices
    \verb|v1| y \verb|v2| en el grafo \verb|g|.
\end{enumerate*}


% B. Método de Pólya para la resolución de problemas
\chapter{Método de Pólya para la resolución de problemas}
\label{Polya}

\section{Método de Pólya para la resolución de problemas matemáticos}

\noindent Para resolver un problema se necesita:
\subsubsection*{Paso 1: Entender el problema}
\begin{itemize}
\item ¿Cuál es la incógnita?, ¿Cuáles son los datos?
\item ¿Cuál es la condición? ¿Es la condición suficiente para determinar la
  incógnita? ¿Es insuficiente? ¿Redundante? ¿Contradictoria? 
\end{itemize}

\subsubsection*{Paso 2: Configurar un plan}
\begin{itemize}
\item ¿Te has encontrado con un problema semejante? ¿O has visto el mismo
  problema planteado en forma ligeramente diferente? 
\item ¿Conoces algún problema relacionado con éste? ¿Conoces algún teorema que
  te pueda ser útil? Mira atentamente la incógnita y trata de recordar un
  problema que sea familiar y que tenga la misma incógnita o una incógnita
  similar. 
\item He aquí un problema relacionado al tuyo y que ya has resuelto ya. ¿Puedes
  utilizarlo? ¿Puedes utilizar su resultado? ¿Puedes emplear su método? ¿Te
  hace falta introducir algún elemento auxiliar a fin de poder utilizarlo? 
\item ¿Puedes enunciar al problema de otra forma? ¿Puedes plantearlo en forma
  diferente nuevamente? Recurre a las definiciones.
\item Si no puedes resolver el problema propuesto, trata de resolver primero
  algún problema similar. ¿Puedes imaginarte un problema análogo un tanto más
  accesible? ¿Un problema más general? ¿Un problema más particular? ¿Un problema
  análogo? ¿Puede resolver una parte del problema? Considera sólo una parte de
  la condición; descarta la otra parte; ¿en qué medida la incógnita queda ahora
  determinada? ¿En qué forma puede variar? ¿Puedes deducir algún
  elemento útil de los datos? ¿Puedes pensar en algunos otros datos apropiados
  para determinar la incógnita? ¿Puedes cambiar la incógnita? ¿Puedes cambiar la
  incógnita o los datos, o ambos si es necesario, de tal forma que estén más
  cercanos entre sí? 
\item ¿Has empleado todos los datos? ¿Has empleado toda la condición? ¿Has
  considerado todas las nociones esenciales concernientes al problema? 
\end{itemize}

\subsubsection*{Paso 3: Ejecutar el plan}
\begin{itemize}
\item Al ejercutar tu plan de la solución, comprueba cada uno de los pasos
\item ¿Puedes ver claramente que el paso es correcto? ¿Puedes demostrarlo? 
\end{itemize}

\subsubsection*{Paso 4: Examinar la solución obtenida}
\begin{itemize}
\item ¿Puedes verificar el resultado? ¿Puedes el razonamiento?
\item ¿Puedes obtener el resultado en forma diferente? ¿Puedes verlo de golpe?
  ¿Puedes emplear el resultado o el método en algún otro problema? 
\end{itemize}

\noindent
\textit{G. Polya ``Cómo plantear y resolver problemas'' (Ed. Trillas, 1978)
  p. 19} 

\section{Método de Pólya para resolver problemas de programación}

\noindent Para resolver un problema se necesita:
\subsubsection*{Paso 1: Entender el problema}
\begin{itemize}
\item ¿Cuáles son las \emph{argumentos}? ¿Cuál es el \emph{resultado}? ¿Cuál es
  \emph{nombre} de la función? ¿Cuál es su \emph{tipo}?
\item ¿Cuál es la \emph{especificación} del problema? ¿Puede satisfacerse la
  especificación? ¿Es insuficiente? ¿Redundante? ¿Contradictoria? ¿Qué
  restricciones se suponen sobre los argumentos y el resultado?
\item ¿Puedes descomponer el problema en partes? Puede ser útil dibujar
  diagramas con ejemplos de argumentos y resultados.
\end{itemize}

\subsubsection*{Paso 2: Diseñar el programa}
\begin{itemize}
\item ¿Te has encontrado con un problema semejante? ¿O has visto el mismo
  problema planteado en forma ligeramente diferente? 
\item ¿Conoces algún problema \emph{relacionado} con éste? ¿Conoces alguna
  función que te pueda ser útil? Mira atentamente el tipo y trata de recordar un
  problema que sea familiar y que tenga el mismo tipo o un tipo similar. 
\item ¿Conoces algún problema familiar con una \emph{especificación} similar? 
\item He aquí un problema \emph{relacionado} al tuyo y que ya has
  resuelto. ¿Puedes utilizarlo? ¿Puedes utilizar su resultado? ¿Puedes emplear
  su método? ¿Te hace falta introducir alguna función auxiliar a fin de poder
  utilizarlo?  
\item Si no puedes resolver el problema propuesto, trata de resolver primero
  algún problema similar. ¿Puedes imaginarte un problema análogo un
  tanto más \emph{accesible}? ¿Un problema más \emph{general}? ¿Un problema más
  \emph{particular}? ¿Un problema \emph{análogo}? 
\item ¿Puede resolver una \emph{parte} del problema? ¿Puedes deducir algún
  elemento útil de los datos? ¿Puedes pensar en algunos otros datos apropiados
  para determinar la incógnita? ¿Puedes cambiar la incógnita? ¿Puedes cambiar la
  incógnita o los datos, o ambos si es necesario, de tal forma que estén más
  cercanos entre sí? 
\item ¿Has empleado todos los datos? ¿Has empleado todas las restricciones
  sobre los datos? ¿Has considerado todas los requisitos de la especificación?
\end{itemize}

\subsubsection*{Paso 3: Escribir el programa}
\begin{itemize}
\item Al escribir el programa, comprueba cada uno de los pasos y funciones
  auxiliares. 
\item ¿Puedes ver claramente que cada paso o función auxiliar es correcta?
\item Puedes escribir el programa en \emph{etapas}. Piensas en los diferentes
  \emph{casos} en los que se divide el problema; en particular, piensas en los
  diferentes casos para los datos. Puedes pensar en el cálculo de los casos
  independientemente y \emph{unirlos} para obtener el resultado final
\item Puedes pensar en la solución del problema descomponiéndolo en problemas
  con datos más simples y uniendo las soluciones parciales para obtener la
  solución del problema; esto es, por \emph{recursión}.
\item En su diseño se puede usar problemas más generales o más
  particulares. Escribe las soluciones de estos problemas; ellas puede servir
  como guía para la solución del problema original, o se pueden usar en su
  solución. 
\item ¿Puedes apoyarte en otros problemas que has resuelto? ¿Pueden usarse?
  ¿Pueden modificarse? ¿Pueden guiar la solución del problema original?
\end{itemize}

\subsubsection*{Paso 4: Examinar la solución obtenida}
\begin{itemize}
\item ¿Puedes comprobar el funcionamiento del programa sobre una colección de
  argumentos? 
\item ¿Puedes comprobar propiedades del programa?
\item ¿Puedes escribir el programa en una forma diferente?
\item ¿Puedes emplear el programa o el método en algún otro programa? 
\end{itemize}

\noindent
Simon Thompson 
\href{http://www.cs.kent.ac.uk/people/staff/sjt/Haskell_craft/HowToProgIt.html}
     {\emph{How to program it}}, 
basado en G. Polya \emph{Cómo plantear y resolver problemas}.


%%%%%%%%%%%%%%%%%%%%%%%%%%%%%%%%%%%%%%%%%%%%%%%%%%%%%%%%%%%%%%%%%%%%%%%%%%%%%%%
%%  Bibliografía                                                            %%
%%%%%%%%%%%%%%%%%%%%%%%%%%%%%%%%%%%%%%%%%%%%%%%%%%%%%%%%%%%%%%%%%%%%%%%%%%%%%%%

% \nocite{Alonso-12a}
\nocite{Alonso-12b}
\nocite{Bird-99a}
\nocite{Cunningham-10a}
\nocite{Daume-06}
\nocite{Davie-92a}
\nocite{Doets-04a}
\nocite{Fokker-96}
\nocite{Hudak-00a}
\nocite{Hudak-12a}
\nocite{Hutton-07a}
\nocite{OSullivan-08a}
\nocite{Rabhi-99a}
\nocite{Polya-65a}
\nocite{Ruiz-04}
\nocite{Thompson-11a}

\addcontentsline{toc}{chapter}{Bibliografía}
\bibliographystyle{abbrv}
\bibliography{Examenes_de_PF_con_Haskell_Vol10}

\end{document}

%%% Local Variables: 
%%% mode: latex
%%% TeX-master: t
%%% End: 

